% Fixed133 independent substantive insertion.
% Exact marked 3-sheet Reidemeister-Schreier presentation of
% H(2), integral Fox/Shapiro augmentation, primitive pro-2
% Nielsen/Tietze eliminations, the rank-five quadratic cup matrix,
% and the complete coefficient-natural degree-two Fox chain split.

\subsection{Actual pro-$2$ Schreier relations and an integral
five-generator Fox reduction in the index-three dyadic field}
\label{subsec:dyadic-f133-actual-syzygy}

The fixed132 zero-sum Schreier exchange identities
are coefficient-level consequences of finite
tame conjugacy.  They are \emph{not}
the relations of $G_F(2)$.
We now construct the relations themselves
by applying marked Reidemeister--Schreier
rewriting to the genuine Roe--Turturean
presentation of $G_{\mathbf Q_2}$.
For the specific index-three field
$F=\mathbf Q_2(\sqrt[3]2)$,
the result is a complete finite
\emph{pro-$2$ group presentation}.
Two successive integral unit-Fox
eliminations reduce its six indexed
relations to one rank-five Demu\v{s}kin
relation, with explicit quadratic
Magnus symbol.  The same transformations
split the finite coefficient
Fox complex \emph{in degree two},
so in this chosen non-pro-$2$
residual sector the former
Schreier relation-syzygy gap
is closed up to a
marked, finite-jet-computable
change of minimal generators.

Throughout put
\[
\Gamma=G_{\mathbf Q_2},\quad
H=G_F,\quad
P_F=H(2)=G_F(2).
\]
Let $\sigma,\tau,x_0,x_1$
be the marked Roe--Turturean
generators as in
Theorem~\ref{thm:dyadic-q2-marked-word-jet-solver}.
The genuine quotient
$\Gamma\twoheadrightarrow S_3$
sends $\sigma$ to a
transposition, $\tau$ to
a three-cycle and
$x_0,x_1$ to the identity.
The right coset representatives
are $t_i=\tau^i$ for
$i\in\{0,1,2\}$.

\begin{theorem}[Marked three-sheet Schreier relators]
\label{thm:dyadic-f133-schreier-pro2-presentation}
In the Schreier presentation
before eliminating its tame
direction, use the ten
pro-$2$ source generators
\begin{equation}\label{eq:dyadic-f133-ten-schreier-generators}
\boxed{
\begin{gathered}
s_0=\sigma,\quad
s_1=\tau\sigma\tau^{-2},
\quad
s_2=\tau^2\sigma\tau^{-1},
\quad t=\tau^3,\\
w_{ij}=\tau^i x_j\tau^{-i}
\quad(i=0,1,2;\ j=0,1).
\end{gathered}}
\end{equation}
The three \emph{literal}
Reidemeister--Schreier
rewrites of the tame
relator
$r_t=\sigma^{-1}\tau\sigma\tau^{-2}$
are
\begin{equation}\label{eq:dyadic-f133-three-tame-relators}
\boxed{
r_{t,0}=s_0^{-1}s_1,\quad
r_{t,1}=s_2^{-1}t s_0t^{-1},
\quad
r_{t,2}=s_1^{-1}s_2t^{-1}.
}
\end{equation}
In every pro-$2$
target they imply
\begin{equation}\label{eq:dyadic-f133-tame-elimination}
\boxed{\qquad
s_1=s_2=s_0=:s,\qquad t=1.
\qquad}
\end{equation}
More precisely the
first and third relations
give $s_1=s_0$ and
$s_2=s_0t$, and
the second then becomes
$s^{-1}ts=t^2$.
No nontrivial element
of a finite $2$-group
can be conjugate to its
square, so $t=1$ in
the universal pro-$2$
quotient.

After this \emph{exact}
tame elimination, put
\[
a_i=w_{i0},\qquad
v_i=w_{i1},\qquad
i\in\mathbf Z/3.
\]
The following are genuine
pro-$2$ group words,
with $1/3\in\mathbf Z_2$
acting by continuous
pro-$2$ powering
\emph{after} the tame
Tietze elimination
$t=1$:
\begin{equation}\label{eq:dyadic-f133-u-and-d-words}
\boxed{
\begin{aligned}
u_{j,i}
 &=\bigl(w_{i,j}w_{i+1,j}
              w_{i+2,j}\bigr)^{1/3},
 & d_i&=u_{0,i}a_i^{-1},\\
z_i&=s^{-1}a_{2i}s,
 & c_i&=[d_i,z_i],\\
g&=s^2,&
D_i&=g^{-1}d_i g,\\
h_i&=[D_i,d_i],
 &H_i&=(g^{-1}a_i g)\,
       a_iD_i d_i^3h_i.
\end{aligned}}
\end{equation}
The indexed wild
Roe--Turturean relators
are the three explicit
pro-$2$ words
\begin{equation}\label{eq:dyadic-f133-three-wild-relators}
\boxed{\quad
R_i
=H_i\,u_{1,i}^{-1}\,
s^{-1}v_{2i}s\,c_i,
\qquad i=0,1,2
\quad(\text{indices mod }3).
\quad}
\end{equation}
They give the exact
seven-generator,
three-relator pro-$2$
presentation
\begin{equation}\label{eq:dyadic-f133-seven-three-presentation}
\boxed{\quad
P_F\cong
\Bigl\langle
s,a_0,a_1,a_2,v_0,v_1,v_2
\ \Bigm|\
R_0,R_1,R_2
\Bigr\rangle_{\mathrm{pro}-2}.
\quad}
\end{equation}
This is a \emph{group-level}
presentation of the actual
maximal pro-$2$ Galois
group, not a coefficient
exchange identity.
\end{theorem}

\begin{proof}
The Schreier formula
for the chosen right
coset representatives is
$t_i g=h_{i,g}t_{j(i,g)}$.
Before using the tame
relation, the three
$\sigma$ edges are
$s_0,s_1,s_2$;
the third $\tau$ edge
is $t=\tau^3$;
the two earlier $\tau$
edges belong to a
spanning tree and
give no generator.
Each wild $x_j$ yields
the three loops $w_{ij}$.
This gives ten
Schreier generators.
Applying the
Schreier word scan to
$r_t=\sigma^{-1}\tau
\sigma\tau^{-2}$ from
initial cosets $0,1,2$
gives exactly
\eqref{eq:dyadic-f133-three-tame-relators}.
The stated elimination
follows by the
element-order argument.

For the wild words, note
the exact group identity
\begin{equation}\label{eq:dyadic-f133-threefold-wild-cycle}
\boxed{\quad
\tau^i(x_j\tau)^3\tau^{-i}
=\widetilde w_{i,j}\widetilde w_{i+1,j}
 \widetilde w_{i+2,j}\tau^3,
\quad
\widetilde w_{k,j}:=
\tau^k x_j\tau^{-k}\ (k\in\mathbf Z).
\quad}
\end{equation}
This is an exact
identity \emph{before}
taking the pro-$2$
quotient.  The
extended indices
obey
$\widetilde w_{k+3,j}
=t\widetilde w_{k,j}t^{-1}$,
where $t=\tau^3$.
Only \emph{after}
the tame relation
forces $t=1$
may all indices
be read modulo three,
giving the displayed
cyclic packets in
\eqref{eq:dyadic-f133-u-and-d-words}.
The right-hand side
belongs to $H$.
The continuous
$\omega_2$ power
$u_j=(x_j\tau)^{\omega_2}$
belongs to the
marked wild
pro-$2$ normal
subgroup.
Its third power
differs from
$(x_j\tau)^3$
by a pro-odd
element in the
same procyclic
closure.  This
pro-odd component
dies in $P_F$.
After $t=\tau^3=1$,
the unique pro-$2$
cube root gives
precisely $u_{j,i}$
in
\eqref{eq:dyadic-f133-u-and-d-words}.
The remaining
wild auxiliary
words are products,
inverses,
conjugates and
commutators.
The Schreier
coset action of
$\sigma$ is
$i\mapsto2i$;
after tame
elimination its
edge value is $s$.
Thus
$\tau^i x_0^\sigma\tau^{-i}$
rewrites as
$s^{-1}a_{2i}s$,
and
$\tau^i x_1^\sigma\tau^{-i}$
as $s^{-1}v_{2i}s$.
The word $\sigma^2$
has the common
Schreier value $s^2$.
Substituting these
expressions into
the actual marked
wild relation
$h_0u_1^{-1}x_1^\sigma c_0$
gives the
three words $R_i$
\eqref{eq:dyadic-f133-three-wild-relators}.

For clarity, the
\emph{marked pro-$2$
condition} has not
been discarded
silently.  A
finite quotient
$H\to Q$ with
$Q$ a $2$-group
induces the usual
finite wreath
representation
$\Gamma\to Q^3\rtimes S_3$:
its three coordinates
are the chosen
Schreier cosets.
The images of
$x_0,x_1$ lie in
the base $Q^3$,
whose normal
closure is a
$2$-group.
Thus every
such wreath
quotient is
admissible for
the marked
Roe--Turturean
presentation
\cite{RoeTurtureanGQ2}.
Conversely a solution
of the indexed
Schreier relators
in an arbitrary
finite $2$-group
produces this
admissible finite
wreath quotient.
The two finite
quotient functors
therefore coincide.
Passing to their
inverse limit proves
the claimed
\emph{exact}
pro-$2$ presentation
\eqref{eq:dyadic-f133-seven-three-presentation}.
This argument uses
the actual marked
source, not the
generally false
assumption that
its two naked
word relations
alone present
the full absolute
Galois group.
\end{proof}

\begin{theorem}[Exact augmented Fox row and a unit pro-$2$ Tietze minor]
\label{thm:dyadic-f133-primitive-wild-fox-minor}
Order the seven
post-tame generators
as
$(s,a_0,a_1,a_2,v_0,v_1,v_2)$.
For the three
\emph{actual}
pro-$2$ relators
$R_i$, the
augmented completed
Fox Jacobian is
\begin{equation}\label{eq:dyadic-f133-augmented-three-seven}
\boxed{
\varepsilon\!\left(
\frac{\partial R_i}
     {\partial s}\right)=0,
\quad
\varepsilon\!\left(
\frac{\partial R_i}
     {\partial a_k}\right)
=\frac43-2\delta_{ik},
\quad
\varepsilon\!\left(
\frac{\partial R_i}
     {\partial v_k}\right)
=\delta_{k,\,2i}-\frac13.
}
\end{equation}
In particular
\begin{equation}\label{eq:dyadic-f133-wild-modtwo-relations}
\boxed{
\begin{aligned}
\bar R_0&=\bar v_1+\bar v_2,\\
\bar R_1&=\bar v_0+\bar v_1,\\
\bar R_2&=\bar v_0+\bar v_2
\end{aligned}
\quad
\text{in }F_7/\Phi(F_7)
\cong\mathbf F_2^7.
}
\end{equation}
The $2\times2$
Fox submatrix of
$(R_1,R_2)$ against
$(v_0,v_1)$ has
the exact unit
determinant
\begin{equation}\label{eq:dyadic-f133-wild-nielsen-unit}
\boxed{
\det
\begin{pmatrix}
-1/3&-1/3\\
-1/3&2/3
\end{pmatrix}
=-\frac13
\in\mathbf Z_2^\times.
}
\end{equation}

Consequently the
assignment of
free pro-$2$
generators
\begin{equation}\label{eq:dyadic-f133-nielsen-basis-change}
\boxed{
\begin{aligned}
&s\mapsto s,\qquad
a_i\mapsto a_i\ (i=0,1,2),\qquad
v_2\mapsto v_2,\\
&v_0\mapsto R_1,\qquad
v_1\mapsto R_2
\end{aligned}}
\end{equation}
defines an
\emph{automorphism}
$\alpha$ of the
free pro-$2$
group $F_7$.
Let $\operatorname{ev}_{0}$
mean substitution
$v_0=v_1=1$
in the \emph{new}
generator coordinates.
Then the five-generator
relator
\begin{equation}\label{eq:dyadic-f133-one-relator-algorithm}
\boxed{\qquad
\mathscr R_F:=
\operatorname{ev}_{0}
\Bigl(\alpha^{-1}(R_0R_1R_2)\Bigr)
\in\widehat F_{2}(s,a_0,a_1,a_2,v_2)
\qquad}
\end{equation}
gives the exact
\emph{minimal}
pro-$2$ presentation
\begin{equation}\label{eq:dyadic-f133-five-one-presentation}
\boxed{\quad
P_F\cong
\langle s,a_0,a_1,a_2,v_2
\mid \mathscr R_F\rangle_{\mathrm{pro}-2}.
\quad}
\end{equation}
The inverse
$\alpha^{-1}$
and $\mathscr R_F$
are recursively
computable in
every finite
Zassenhaus/Frattini
quotient of the
free pro-$2$
group, using
only the three
explicit indexed
words and
continuous $1/3$
powering.
No closed finite
abstract-word formula
for all its higher
pro-$2$ digits
is claimed.
\end{theorem}

\begin{proof}
The abelian
exponent sums
of $u_{j,i}$
are one third
the sum of
the three
corresponding
wild generators.
All commutators
and conjugations
have the expected
zero or unchanged
exponent sums.
In particular
\[
[H_i]_{\mathrm{ab}}
=2[a_i]+4([u_{0,i}]-[a_i])
=\frac43\sum_k[a_k]-2[a_i],
\]
while
$[u_{1,i}^{-1}
s^{-1}v_{2i}s]_{\mathrm{ab}}
=-\frac13\sum_k[v_k]+[v_{2i}]$.
This proves
\eqref{eq:dyadic-f133-augmented-three-seven}.
Modulo two,
$4/3$ and $-2$
vanish while
$-1/3=1$.
This yields
\eqref{eq:dyadic-f133-wild-modtwo-relations}
and the displayed
unit minor.

The images of
the seven new
words in
$F_7/\Phi(F_7)$
are a basis
because the
two-by-two
block in
\eqref{eq:dyadic-f133-wild-nielsen-unit}
is invertible
modulo two,
while the
other five
generators are
unchanged.
The pro-$2$
Burnside basis
theorem therefore
makes
\eqref{eq:dyadic-f133-nielsen-basis-change}
a surjective
endomorphism
of the finitely
generated free
pro-$2$ group;
the Hopfian
property makes
it an automorphism.
Under its inverse,
the first two
relators $R_1,R_2$
become the new
generators
$v_0,v_1$.
Killing these
relators eliminates
those generators
integrally and
leaves the image
of $R_0$ as
the only relation.
Multiplying by
$R_1R_2$ before
the substitution
does not change
the normal
closure, and
makes the
quadratic Magnus
calculation
transparent.
This gives
\eqref{eq:dyadic-f133-five-one-presentation}.

For each fixed
finite pro-$2$
quotient of
the source free
group, exponent
$1/3$ is a
well-defined
finite power
with exponent
inverse to three
modulo the
relevant group
exponent.
The induced
basis automorphism
is bijective
on that finite
quotient, and
its inverse
can be computed
by finite
enumeration or
successive
Frattini lifts.
Taking compatible
quotients
determines
$\alpha^{-1}$
in the inverse
limit, giving
the stated
finite-stage
effectivity.
\end{proof}

\begin{theorem}[The surviving marked quadratic Fox--Demu\v{s}kin symbol]
\label{thm:dyadic-f133-explicit-five-cup-symbol}
Write $S,A_0,A_1,A_2,V$
for the mod-$2$
Magnus variables
corresponding to
the five generators
$(s,a_0,a_1,a_2,v_2)$.
The unique remaining
relator
$\mathscr R_F$
has the exact
integral augmented
linear data
\begin{equation}\label{eq:dyadic-f133-minimal-linear-fox}
\boxed{
\varepsilon(\partial_s\mathscr R_F)=0,\qquad
\varepsilon(\partial_{a_i}\mathscr R_F)=2
\quad(i=0,1,2),\qquad
\varepsilon(\partial_{v_2}\mathscr R_F)=0.
}
\end{equation}
Its degree-two
mod-$2$ Magnus
symbol is
\begin{equation}\label{eq:dyadic-f133-minimal-quadratic-fox}
\boxed{\qquad
\sigma_2(\mathscr R_F)
=A_0^2+A_1^2+A_2^2+SV+VS.
\qquad}
\end{equation}
The resulting
cohomological
cup matrix is
\begin{equation}\label{eq:dyadic-f133-cup-rank-five}
\boxed{
B_F=
\begin{pmatrix}
0&0&0&0&1\\
0&1&0&0&0\\
0&0&1&0&0\\
0&0&0&1&0\\
1&0&0&0&0
\end{pmatrix},
\qquad
\det_{\mathbf F_2}B_F=1.
}
\end{equation}
In the
$\mathbf F_2$
cohomology basis
\begin{equation}\label{eq:dyadic-f133-demushkin-cup-change}
\boxed{
a_\star=a_0+a_1+a_2,\quad
b_\star=a_0+a_1,\quad
c_\star=a_1+a_2,\quad
s,\quad v_2,
}
\end{equation}
the symmetric
cup form is
$\langle1\rangle
\oplus H\oplus H$,
where
$H=\begin{psmallmatrix}0&1\\1&0\end{psmallmatrix}$.
Thus this is
a nondegenerate
\emph{nonalternating}
rank-five dyadic
Demu\v{s}kin form.
The integral
first Fox ideal
is precisely $(2)$,
so the five-generator
relation has
\[
P_F^{\mathrm{ab}}
\cong\mathbf Z_2^{\,4}
\oplus\mathbf Z/2\mathbf Z.
\]
These two
independent
calculations
match the known
$[F:\mathbf Q_2]+2=5$,
$q(F)=2$,
full-orientation
local Demu\v{s}kin
classification, but
do \emph{not}
rely on that
classification
to establish
nondegeneracy.
\end{theorem}

\begin{proof}
Adding the exact
augmented first
Fox rows from
\eqref{eq:dyadic-f133-augmented-three-seven}
gives
\[
\varepsilon(\partial(R_0R_1R_2))
=(0,2,2,2,0,0,0).
\]
The entries in
the two subsequently
eliminated variables
$v_0,v_1$ vanish
\emph{integrally}.
Consequently the
pro-$2$ Nielsen
elimination does
not alter these
linear values,
which proves
\eqref{eq:dyadic-f133-minimal-linear-fox}.

For the quadratic
symbol, work in
the truncated
noncommutative
Magnus algebra
$\mathbf F_2
\langle\!\langle
S,A_0,A_1,A_2,
V_0,V_1,V_2
\rangle\!\rangle/
(\mathrm{degree}\ge3)$.
For a pro-$2$
word with expansion
$1+L+Q+O(3)$,
continuous cubic
rooting satisfies
\begin{equation}\label{eq:dyadic-f133-magnus-one-third}
\boxed{\quad
(1+L+Q)^{1/3}
=1+L+Q+L^2
\pmod{(\deg\ge3),2}.
\quad}
\end{equation}
Indeed $\binom{1/3}{1}$
and $\binom{1/3}{2}=-1/9$
are both odd
$2$-adic integers.
Use this identity,
the ordered
commutator convention
and the indexed
word formulas
\eqref{eq:dyadic-f133-u-and-d-words}
and
\eqref{eq:dyadic-f133-three-wild-relators}
to expand
$R_0R_1R_2$
through Magnus
degree two.
Its linear
coefficients in
all three
$V_i$ are
\emph{exactly}
zero; therefore
replacing
$V_0,V_1,V_2$
by their
common degree-one
class $V$
cannot change
its quadratic
term even when
the eliminated
group generators
have higher-order
nonlinear
corrections.
The resulting
quadratic expression
is exactly
\[
A_0^2+A_1^2+A_2^2+SV+VS,
\]
which proves
\eqref{eq:dyadic-f133-minimal-quadratic-fox}.
It is also
reproduced by the
standalone finite
Magnus/Fox
verification
attached to
this version.

The ordered
quadratic Fox--cup
comparison
of
Theorem~\ref{thm:quadratic-fox-cup}
gives the matrix
\eqref{eq:dyadic-f133-cup-rank-five}.
Its determinant
is one.
The three
$A_i$ squares
give a rank-three
identity form,
while $(S,V)$
form an
independent
hyperbolic pair.
In the proposed
new basis
$a_\star$
has self-pairing
one and is
orthogonal to
$b_\star,c_\star$;
those two
have zero squares
and pair to one.
The final pair
$(s,v_2)$ is
hyperbolic and
orthogonal to
all three.
This is the
displayed
$\langle1\rangle
\oplus H\oplus H$
normal form.
Finally the
linear exponent
row is $(0,2,2,2,0)$,
whose Smith ideal
is $(2)$, giving
the abelianization
formula.
\end{proof}

\begin{theorem}[Integral three-sheet Fox--Shapiro matrix and five unit pivots]
\label{thm:dyadic-f133-cover-fox-shapiro-matrix}
The original
marked two-relator
presentation has a
three-sheet
Schreier Fox
$2$-skeleton
with three vertices,
twelve generating
edges and six
indexed relation
faces.  Let
$T=\mathbf Z_2$
be the trivial
$H$-coefficient
lattice.
Via
coinduction,
\[
V=\operatorname{Coind}_H^\Gamma T
\cong\mathbf Z_2^3
\]
has the marked
generator actions
\begin{equation}\label{eq:dyadic-f133-three-sheet-actions}
\boxed{
\Sigma=
\begin{pmatrix}
1&0&0\\0&0&1\\0&1&0
\end{pmatrix},
\qquad
\Upsilon=
\begin{pmatrix}
0&1&0\\0&0&1\\1&0&0
\end{pmatrix},
\qquad
\rho_V(x_0)=\rho_V(x_1)=I_3.
}
\end{equation}
Put $N=I+\Upsilon+\Upsilon^2$
and $P=N/3$.
In coordinates $(a,b,c,d)\in V^4$
for the four marked
generator cochains,
the complete
coefficient Fox
differentiels are
\begin{equation}\label{eq:dyadic-f133-coinduced-full-fox}
\boxed{
\begin{aligned}
d^0(v)
 &=\bigl((\Sigma-I)v,
       (\Upsilon-I)v,0,0\bigr),\\
d^1(a,b,c,d)
 &=\Bigl(
 \Sigma^{-1}(\Upsilon-I)a+
  (\Sigma^{-1}-I-\Upsilon)b,\\
 &\hspace{2.9em}
 3Pb+(4P-2I)c+
      (\Sigma^{-1}-P)d
 \Bigr).
\end{aligned}}
\end{equation}
The two $\tau$
spanning-tree edges,
corresponding to
cochain coordinates
$b_0,b_1$, have
unit vertex--edge
boundary.
Contract them,
and order the ten
remaining edge
cochains as
\[
(s_0,s_1,s_2,t,
a_0,a_1,a_2,v_0,v_1,v_2).
\]
After multiplying
the three wild
face rows by
the $2$-adic unit
$3$, the exact
augmented integral
Fox matrix is
\begin{equation}\label{eq:dyadic-f133-six-ten-fox-matrix}
\boxed{
\mathsf J^{(3)}_{6,10}=
\begin{pmatrix}
-1&1&0&0&0&0&0&0&0&0\\
1&0&-1&0&0&0&0&0&0&0\\
0&-1&1&-1&0&0&0&0&0&0\\ \hline
0&0&0&3&-2&4&4&2&-1&-1\\
0&0&0&3&4&-2&4&-1&-1&2\\
0&0&0&3&4&4&-2&-1&2&-1
\end{pmatrix}.
}
\end{equation}
Its mod-two rank
is exactly five.
Two explicit
integral minors,
using one-based
column indices,
are
\begin{equation}\label{eq:dyadic-f133-fox-minor-certificates}
\boxed{
\begin{aligned}
\det\bigl(
\mathsf J^{(3)}_{\{1,\dots,5\},
 \{1,2,4,8,9\}}\bigr)&=-3,\\
\det\bigl(
\mathsf J^{(3)}_{\{1,\dots,6\},
 \{1,2,4,5,8,9\}}\bigr)&=-18.
\end{aligned}}
\end{equation}
Therefore the
exact $2$-adic
Smith factors of
the augmented
face differential
are
\begin{equation}\label{eq:dyadic-f133-fox-six-ten-smith}
\boxed{\quad
(1,1,1,1,1,2).
\quad}
\end{equation}
After the two
spanning-tree
contractions,
the integral
cohomology of
this finite
three-sheet
coefficient complex
is
\begin{equation}\label{eq:dyadic-f133-h-trivial-ledger}
\boxed{\quad
H^0(H,\mathbf Z_2)=\mathbf Z_2,\quad
H^1(H,\mathbf Z_2)=\mathbf Z_2^4,\quad
H^2(H,\mathbf Z_2)=\mathbf F_2.
\quad}
\end{equation}
The first five
unit pivots are
the precise
linear shadow of
the three tame
and two wild
group-level
Tietze eliminations
above.
\end{theorem}

\begin{proof}
Coinduction
identifies the
three coset fibers
with the three
coordinates of
$V$.  The Schreier
table
\eqref{eq:dyadic-s3-schreier-table}
gives the
$\Sigma,\Upsilon$
matrices.
On the
affine semidirect
extension
$V\rtimes S_3$,
one has
$\Sigma^2=I$,
$\Upsilon^3=I$,
$\Sigma^{-1}\Upsilon\Sigma
=\Upsilon^2$.
For an arbitrary
translation $b+c$,
the $\omega_2$
projection of
$(\Upsilon,b+c)$
is the pure
translation
$(I,P(b+c))$,
since the
cube is
$(I,N(b+c))$
and three is
invertible.
Thus, in the
literal marked
wild relator,
the affine
auxiliary translation
is
$d_0=Pb+(P-I)c$.
The remaining
commutators of
pure translations
vanish.
The wild
relation evaluates
to
\[
2c+4d_0-P(b+d)
+\Sigma^{-1}d
=
3Pb+(4P-2I)c+
(\Sigma^{-1}-P)d.
\]
The tame relation
evaluates to the
first component of
\eqref{eq:dyadic-f133-coinduced-full-fox}.
This gives the
full completed
Fox differential
by the crossed
derivation--Fox
identity.

The first
two $\tau$ edge
values of
$(\Upsilon-I)v$
are $v_1-v_0$
and $v_2-v_1$,
which form
a primitive
rank-two vertex
coboundary.
Use them to
cancel the
two tree edges.
Set their
cochain values
to zero in
the chosen
strict gauge,
and retain the
third $\tau$
edge as the
$t=\tau^3$
generator.
Applying the
row-unit change
which multiplies
the three wild
faces by three
gives exactly
\eqref{eq:dyadic-f133-six-ten-fox-matrix}.
Its first
three rows
are the augmented
Fox rows of
\eqref{eq:dyadic-f133-three-tame-relators};
its last three
rows are the
indexed wild
rows before
tame elimination.

Direct
mod-two row
reduction gives
rank five,
with pivot
columns
$\{1,2,4,8,9\}$.
The first
displayed minor
is an odd
unit.
Every
$6\times6$
minor is
even because
the mod-two
rank is five;
the second
displayed
minor has
valuation exactly
one.  The Smith
invariants are
therefore five
units and one
factor two.
The remaining
degree-zero
and degree-one
ranks give
$H^0=\mathbf Z_2$,
$H^1=\mathbf Z_2^4$,
and $H^2=\mathbf F_2$,
also agreeing
with local
class field
theory for the
degree-three
dyadic field.
\end{proof}

\begin{theorem}[An integral Fox--Shapiro chain equivalence through degree two]
\label{thm:dyadic-f133-integral-schreier-fox-chain}
Let
\[
\Lambda_F=\mathbf Z_2[[P_F]]
\]
be the completed
pro-$2$ group algebra,
and let $M$ be
\emph{any} finite
projective complete
coefficient module
whose $H$ action
factors continuously
through $P_F$.
Let
$C_{\rm cov}^\bullet(H,M)$
be the three-sheet
marked
Reidemeister--Schreier
Fox $2$-skeleton:
three vertex copies
of $M$,
twelve generating-edge
copies, and six
indexed relation-face
copies, with
\emph{all} Fox
derivatives evaluated
in $M$.
Let
\[
C_{\rm min}^\bullet(P_F,M):
\qquad
M\longrightarrow M^5\longrightarrow M
\]
be the standard
completed Fox
complex of
the \emph{actual}
one-relator
pro-$2$
presentation
\eqref{eq:dyadic-f133-five-one-presentation}.
Then there is a
\emph{chosen strict
integral chain
isomorphism}
\begin{equation}\label{eq:dyadic-f133-integral-chain-split}
\boxed{
\begin{aligned}
C_{\rm cov}^\bullet(H,M)
\cong {}&
C_{\rm min}^\bullet(P_F,M)\\
&\oplus
\bigl[M^2\xrightarrow{I}M^2\bigr]_{[0,1]}
\oplus
\bigl[M^5\xrightarrow{I}M^5\bigr]_{[1,2]}.
\end{aligned}}
\end{equation}
The degree-wise
ranks on the
two sides are
exactly $(3,12,6)$.
Every elementary
chain transformation
is integral over
$\Lambda_F$:
the two spanning-tree
pivots, three
tame Nielsen
pivots, and
two wild pro-$2$
Nielsen pivots
are all units.
The changes are
specified by the
Reidemeister--Schreier
word scan,
\eqref{eq:dyadic-f133-three-tame-relators},
\eqref{eq:dyadic-f133-three-wild-relators},
and the pro-$2$
automorphism
\eqref{eq:dyadic-f133-nielsen-basis-change};
their completed
degree-one matrices
are the associated
Fox Jacobians.

The minimal complex
is an \emph{exact}
length-two
projective
resolution after
passing to
$\Lambda_F$,
because $P_F$ is
a rank-five
Demu\v{s}kin group.
Consequently
\begin{equation}\label{eq:dyadic-f133-complete-derived-fact}
\boxed{\quad
C_{\rm cov}^\bullet(H,M)
\simeq R\Gamma(H,M)
\simeq C_{\rm min}^\bullet(P_F,M),
\quad}
\end{equation}
compatibly with
derived coefficient
base change.
In the entire
framed residual
$S_3$ family,
the restricted
universal
coefficient
representation of
$H$ has pro-$2$
image, so the
comparison applies
\emph{uniformly}
to those
coefficient charts,
not just to
the fixed natural
or adjoint
$S_3$ lattice.

Moreover the
change is
\emph{recursively
native-explicit}
at every prescribed
finite pro-$2$
word or coefficient
jet: compute the
inverse of the
seven-generator
basis automorphism
in the finite
source quotient,
evaluate its
completed Fox
derivatives, and
cancel the five
declared unit
relation pivots.
The group-algebra
degree-two syzygy
problem for this
\emph{specific}
$P_F$ is thereby
reduced to the
single genuine
Demu\v{s}kin
relation and
not to an
uncontrolled
three-relator
complex.

The selected
inverse and
strict matrices
depend on the
marked generators,
ordered cosets
and elimination
order.  No
canonical
presentation-independent
comparison is
claimed for
arbitrary dyadic
fields or
arbitrary
non-pro-$2$
residual
quotients.
\end{theorem}

\begin{proof}
Schreier
rewriting is a
literal change
of cellular
coordinates
on the
finite-index
cover of the
marked relation
$2$-complex.
For the three
chosen cosets
it gives
cochain degrees
$(3,12,6)$
with the
specified
indexed edge
and face
coordinates.
The two
tree edges
$\tau:0\to1$
and $\tau:1\to2$
have identity
vertex--edge
coboundaries.
Cancelling those
two contractible
degree-zero--one
pairs leaves
one vertex,
ten edges and
six faces.

The tame
relations
\eqref{eq:dyadic-f133-three-tame-relators}
first replace
$s_1$ by $s_0$,
then $s_2$ by
$s_0t$,
and finally kill
$t$ with the
pro-$2$ relation
$t^s=t^2$.
Each corresponding
augmented Fox
pivot is a
$2$-adic unit;
the actual
completed
Fox pivot
lies in
$\Lambda_F^\times$,
because this
ring is local
with maximal
ideal generated
by two and the
augmentation
ideal.
The corresponding
three relation-face
and generating-edge
pairs therefore
split as identity
degree-one--two
chain summands.
Their remaining
relation words
are precisely
the seven-generator
$R_0,R_1,R_2$.

For the two
wild pivots,
Theorem~\ref{thm:dyadic-f133-primitive-wild-fox-minor}
provides the
\emph{actual pro-$2$
free-group
automorphism}
$\alpha$.
Its completed
Fox Jacobian is
invertible over
the relevant
completed group
algebra by the
Fox chain rule
and the
inverse automorphism.
In the new
free basis the
relations $R_1,R_2$
become two
generator letters.
The resulting
$2\times2$
relation-to-generator
Fox block is
the identity.
After elementary
row operations
it splits as
two further
contractible
degree-one--two
pairs.
The sole
remaining
relation is
$\mathscr R_F$.
All of these
are genuine
pro-$2$ Tietze
operations and
hence induce
literal Fox
chain changes;
none is a
mere numerical
Smith reduction.

The remaining
three-term Fox
complex is the
minimal one-relator
Demu\v{s}kin
resolution.
Its exactness
over $\Lambda_F$
follows from
the finite-rank
pro-$2$
Poincar\'e
duality property
of $P_F$
\cite{LabuteClassification,NSW}.
The two cancelled
tree pairs and
five cancelled
relation--edge
pairs produce
exactly
\eqref{eq:dyadic-f133-integral-chain-split}.
Applying
continuous Hom
to any complete
coefficient
module gives
the cochain
version and
compatibility
with coefficient
base change.
On the residual
$S_3$ chart,
the Sylow
preimage $H$
has residual
image $C_2$;
its universal
framed
congruence kernel
is pro-$2$,
so its full
image is pro-$2$.
Thus the
factorization
through $P_F$
is valid across
that entire
residual sector.

For finite-stage
effectivity,
the Reidemeister--Schreier
rewriting of each
marked word
is finite, and
each continuous
$1/3$ power
is a finite
integer power
in any fixed
finite pro-$2$
quotient.
The inverse
of the pro-$2$
basis automorphism
is computable
in these finite
quotients by
its bijective
finite action.
Completed Fox
differentiation
is continuous
and can be
evaluated at
each chosen
finite jet;
the five
unit pivots
are inverted
by geometric
series in
the complete
local group
algebra.
Hence every
requested
finite-stage
chain matrix
terminates.
The chosen
inverse-limit
maps give the
claimed
strict integral
chain equivalence,
with no assertion
of a universally
finite ordinary
word for
$\mathscr R_F$.
\end{proof}

\begin{corollary}[Closure of the fixed $S_3$ Schreier syzygy frontier]
\label{cor:dyadic-f133-fixed-s3-syzygy-closure}
For
$F=\mathbf Q_2(\sqrt[3]2)$,
the three wild
Schreier relators
in fixed128
are now reduced
by \emph{actual}
pro-$2$ relation
changes to the
single minimal
Demu\v{s}kin relator
$\mathscr R_F$.
The surviving
second Fox row
has a computed
nondegenerate
rank-five
cup symbol,
the exact
integral first
Fox ideal $(2)$,
and a complete
length-two
projective
resolution over
$\mathbf Z_2[[P_F]]$.
For all
continuous
restricted
coefficient
representations
factoring through
$P_F$, the
index-three
Schreier
covering
cochain complex
and the
minimal Fox
complex are
strictly
chain-equivalent
after exactly
seven declared
contractible
cancellations.

This resolves
the \emph{group
relation-syzygy}
problem for the
selected genuine
non-pro-$2$
$S_3$ residual
sector, not
merely the
previous natural
and adjoint
coefficient
examples.
It is a
\emph{choice-dependent}
finite-jet-computable
bridge.
A functorial
minimal Fox
coordinate
system
simultaneously
covering every
non-pro-$2$
residual
quotient and
every dyadic
base field
is a distinct
stronger problem.
\end{corollary}

\begin{proof}
The group
presentation,
the two unit
Nielsen eliminations
and the
one-relator
symbol are
Theorems~\ref{thm:dyadic-f133-schreier-pro2-presentation},
\ref{thm:dyadic-f133-primitive-wild-fox-minor},
and
\ref{thm:dyadic-f133-explicit-five-cup-symbol}.
The coefficient-uniform
group-algebra
Fox chain
equivalence is
Theorem~\ref{thm:dyadic-f133-integral-schreier-fox-chain}.
No additional
independent
relation or
coefficient
hypothesis is
needed.
\end{proof}

\begin{warning}[Marked finite-quotient universality versus global strictness]
\label{warn:dyadic-f133-pro2-marked-scope}
The explicit
presentation is
for the
\emph{maximal
pro-$2$ quotient
of one chosen
index-three
Sylow preimage}
$H=G_F$.
Its proof uses
the actual
marked pro-$2$
normal-closure
condition of
Roe--Turturean,
carried through
finite wreath
quotients.
It is not
the false
unmarked
interpretation
of the two
full-group
word relators.
The resulting
strict
Fox--Shapiro
equivalence is
integral and
valid for
all $H$-coefficient
systems
factoring
through $P_F$,
but its marked
inverse automorphism
and elimination
gauge are
not naturally
independent
of Schreier
choices.
No uniform
closed finite
word formula
over the
entire
family of
dyadic fields,
or preferred
raw Herr
strict
representative,
follows from
this computation.
\end{warning}
